%!TEX root=./neurips_2023.tex
\section{Experiments}
\label{sec:exp}
In this section, we empirically study our approach in five diverse scenarios across different arithmetic operations, and then analyze the effect of PEM initialization and the weight hyperparameter $\lambda$.
% we present the results of our experiments in which we manipulate the capabilities of the PEMs to achieve high flexibility under a training-free setting. 
% We integrate modules trained on different data distributions to facilitate generalization (\textsection \ref{sec:datasetsplit}), transfer knowledge between tasks (\textsection \ref{sec:multitask}), forget certain abilities via negation (\textsection \ref{sec:neg}), and further obtain new skills on an unseen task-domain combination by composing donor modules (\textsection \ref{sec:combine}).
% It is worth noting that the aforementioned operations are performed under the premise that the PEMs have identical pretrained models, module architectures, and module initializations. 
% In addition to PEMs including LoRA~\citep{hu2021lora} and \ia~\citep{liu2022few}, we also conduct experiments on task vectors~\citep{ilharco2022editing}. They are weight difference between full-finetuned model and pretrained one, and this approach served as a reference point rather than a comparable method. 
% To account for mode connectivity~\citep{ainsworth2022git}, in \textsection \ref{sec:analysis} we analyze the effects of divergent initialization on combination performance by merging LoRAs initialized with different random seeds.
\subsection{General Setup}
Throughout the experiments, we fix the pretrained model checkpoints and the architecture of PEMs to be composed the same within each scenario, which are the necessary conditions for arithmetic operations. 
We experiment with LoRA and \ia~ for each scenario unless otherwise specified. We also perform arithmetic operations in the full finetuning (FFT) setting as in~\citet{ilharco2022editing} for a reference point. 
We emphasize that the full finetuning results are not directly comparable to ours since the motivation of this work is composing parameter-efficient modules.
We keep the initialization of the composing PEMs the same for potentially better linear connectivity, while we perform analysis in \textsection\ref{sec:analysis} on the effect of different initialization.
We note that only the $\mA$ matrix in LoRA may be initialized differently -- the $\vl$ vectors in \ia~are all initialized as ones by design as described in \textsection\ref{sec:background}. 
$\lambda$ is the only tunable hyperparameter in our method.
Below for each scenario, we will briefly introduce their setup, and please refer to Appendix~\ref{sec:apx_exp_setup} for complete setup details of all the experiments.
% \jh{To ensure that both modules are inside the same loss minima basin as much as possible after training, identical pretrained models, module architectures, and module initialization are used throughout our experiments to enhance mode connectivity.}

% We vary lambda from 0 to 1 with a step size of 0.02 to optimize the merge setting.

% Full finetuning as a reference point, not a comparable method.

\subsection{Composition for Distribution Generalization}
\label{sec:datasetsplit}

\paragraph{Setup:}
% Training on different distributions enables the PEMs to capture wider range of patterns and insights, resulting in improved robustness and performance.
% Additionally, PEMs can be combined in the weight space with ease, making it an efficient and pragmatic approach.
In this setting, we aim to combine PEMs trained on the same task but divergent distributions, to improve the model's generalization. 
To this end, we follow~\citet{jin2022dataless} to construct a synthetic setting: we select two training subsets from the datasets,
each with imbalanced labels and distinct distributions. Subsequently, we train two separate PEMs on the two subsets respectively and merge them through $\tm = \lambda\ta + (1-\lambda)\tb$. 
We then assess the individual and combined PEMs using the original validation data -- designed to reflect the performance on the union of the subset distributions -- in order to determine whether the merged PEM demonstrates improved generalization capabilities.
%  and merge them for each task with a varied merging proportion, $\lambda$, ranging from $0$ to $1$ with a step size of $0.02$. The merged model, expressed as $\theta_{combined}=\lambda \theta_0 \oplus (1-\lambda) \theta_1$, exhibits superior performance compared to both individual subsets, without any post-processing. 
% \paragraph{Dataset:} 
We work on MNLI~\citep{williams2017broad}, RTE~\citep{giampiccolo2007third}, CoLA~\citep{warstadt2019cola}, SST2~\citep{socher2013recursive}, MRPC~\citep{dolan2005automatically}, QNLI~\citep{rajpurkar2016squad}, QQP~\citep{iyer2017first}, and STS-B~\citep{cer2017semeval} datasets from the GLUE~\citep{wang2018glue} task collections.
Please see Appendix~\ref{sec:apx_exp_setup} on how we construct two distinct subsets from each of the task. 
% We split the datasets by randomly selecting one label and assigning $80\%$ of the label's samples to one subset, and putting the remaining $20\%$ in the other. After that, we randomly distribute samples from the remaining labels to the two subsets to make the number of samples in the subsets equal. For regression task STS-B with values ranging from $0$ to $5$, samples with values above $2.5$ are considered as the selected label.
We adopt RoBERTa-base~\citep{liu2019roberta} as the base model.
The aforementioned datasets are evaluated using accuracy except CoLA, for which we use Matthews Correlation Coefficient (MCC), and STS-B, which we evaluate using the Spearman's rank correlation coefficient.

% \textbf{Models.} 
% We train the encoder-only model RoBERTa-base~\citep{liu2019roberta} under full finetuning, LoRA and \ia  setting. 
% The shared weights of the PEMs, which are theoretically mergable, incorporate task-specific classification heads, as they have the same size.
% Task-specific classification heads with identical sizes are included in the shared weights of the PEMs, which are theoretically mergable.


%%%%%%%%%%%%%%%%%%%%%%%%%%%%%%%%%%%%%%%%
% 8 task left,right,max,improvements


\begin{table*}[]
\caption{The validation results of PEMs trained on both subsets ($s_0$, $s_1$) and merged PEM ($m$). 
``FFT'' represents full finetuning.
% The `$\uparrow$' and numbers with smaller font indicate the average performance improvements. 
We denote the absolute performance change of merged PEM compared to the average results of the two individual PEMs.
We report MCC for CoLA, Spearman's $\rho$ for STS-B, and accuracy for others.
Full-dataset LoRA-tuning results are provided as a reference point, which requires all data in one-way training.
The tuning results for the full dataset using LoRA are provided as a reference point where both subsets of the data are used together for training.
% The datasets are evaluated using percentage accuracy for all except CoLA, for which Matthews Correlation Coefficient (MCC) is used, and STS-B, which evaluated by Spearman rank correlation coefficient (Spearman's $\rho$).
}
\label{table:split}
\centering
\resizebox{1 \columnwidth}{!}{
\begin{tabular}{llllllllll}
\toprule
\multicolumn{2}{l}{Method} & MNLI & RTE & SST-2 & MRPC & QNLI & QQP & CoLA & STS-B \\ \midrule
\multirow{4}{*}{\rotatebox[origin=c]{90}{FFT}} & fullset & 76.6 & 75.8 & 92.5 & 88.5 & 85.9 & 81.8 & 0.56 & 0.90 \\
& $s_0$ & 72.0 & 72.9 & 90.4 & 85.8 & 83.4 & 79.2 & 0.42 & 0.88 \\
 & $s_1$ & 71.9 & 67.5 & 92.0 & 88.5 & 83.2 & 81.5 & 0.52 & 0.89 \\
 & $m$ & 74.2\up{2.3} & 75.1\up{4.9} & 92.1\up{0.9} & 89.2\up{2.1} & 83.8\up{0.5} & 81.9\up{1.5} & 0.55\up{0.07} & 0.89\up{0.01} \\ \midrule
 % & $\Delta$ & 2.3 & 4.9 & 0.9 & 2.1 & 0.5 & 1.5 & 0.071 & 0.005 \\ 
\multirow{4}{*}{\rotatebox[origin=c]{90}{LoRA}} & fullset & 87.1 & 79.8 & 95.0 & 89.2 & 93.4 & 90.2 & 0.63 & 0.91 \\
 & $s_0$ & 71.4 & 72.2 & 92.2 & 86.3 & 83.1 & 79.0 & 0.50 & 0.88 \\
 & $s_1$ & 72.3 & 69.0 & 91.9 & 87.7 & 83.0 & 80.8 & 0.51 & 0.89 \\
 & $m$ & 73.5\up{1.6} & 75.8\up{5.2} & 92.2\up{0.2} & 88.0\up{1.0} & 83.3\up{0.2} & 81.1\up{1.2} & 0.52\up{0.01} & 0.89\up{0.01} \\ \midrule
 % & $\Delta$ & 1.6 & 5.2 & 0.2 & 1.0 & 0.2 & 1.2 & 0.012 & 0.011 \\
\multirow{4}{*}{\rotatebox[origin=c]{90}{\ia}} & fullset & 75.9 & 74.0 & 92.3 & 87.3 & 84.7 & 80.8 & 0.56 & 0.89 \\ 
& $s_0$ & 71.7 & 72.9 & 90.8 & 85.8 & 83.0 & 78.3 & 0.44 & 0.87 \\
 & $s_1$ & 71.7 & 68.2 & 91.2 & 88.0 & 82.5 & 80.8 & 0.50 & 0.90 \\
 & $m$ & 74.0\up{2.3} & 74.7\up{4.0} & 92.3\up{1.3} & 88.2\up{1.3} & 84.8\up{2.0} & 81.3\up{1.8} & 0.50\up{0.03} & 0.90 \up{0.01}\\ \bottomrule
 % & $\Delta$ & 2.3 & 4.0 & 1.3 & 1.3 & 2.0 & 1.8 & 0.031 & 0.013 \\ 
\end{tabular}}
\vspace{-10pt}
\end{table*}

%% fig 曲线图 in appendix
%% exact 比例 in appendix
%% training hyperparameters in appendix

%%%%%%%%%%%%%%%%%%%%%%%%%%%%%%%%%%%%%%%%


\paragraph{Results:}
% In Table \ref{table:split}, we demonstrate the performance of full finetuning models and PEMs trained on the same task but on different distributions, before and after combination. 
We show the results in Table~\ref{table:split}.
After combination, the merged PEM achieves consistent improvement compared to the average performance of two individual PEMs. 
For example, the merged LoRA module and the merged \ia~module obtain gains of 5.2 and 4.0 absolute points respectively on RTE.
% there is a noticeable performance improvement on each task. 
% LoRA and \ia~show up to $1.57\%$ and $2.12\%$ average accuracy increases respectively. 
% \jz{slight change}
Our findings indicate that modular learning permits the integration of abilities via addition. As a consequence, the PEFT approach is capable of not only achieving the same level of performance as full finetuning but also excelling in terms of module composition. This highlights the substantial capabilities of PEFT.
% Furthermore, our results demonstrate that 
% PEMs significantly outperform full finetuning on specific tasks, as evidenced by the $5.23\%$ improvement on the RTE task by LoRa and the $2.03\%$ improvement on the QNLI task by \ia.
% \jz{integrating abilities by addition is available in modular learning. Therefore,}
% PEFT methods can not only match the performance of full finetuning but also performs well concerning module composition, illustrating the significant potential for PEFT.
Analysis of the results change as $\lambda$ varies can be found in Appendix~\ref{sec:apx_lambda}.
% The evaluation results of the merged module, as lambda varies, can be observed in Figure \ref{fig:split0}, which shows that while the RTE results are relatively robust to the $\lambda$ change, the MNLI performance is sensitive to $\lambda$.
% Additionally, the exact distribution of the subsets are listed in Table \ref{table:split0}. The hyperparameters used for full finetuning, LoRA, and \ia tuning can be found in Appendix~\ref{sec:apx_split}.


\subsection{Composition for Multi-Tasking}
\label{sec:multitask}

\paragraph{Setup:}
We examine whether PEMs trained on different tasks could be merged together for multi-task learning. 
Specifically, we follow~\citet{matena2022merging} and select MNLI and RTE as two tasks to be merged.\footnote{We select MNLI and RTE based on the full finetuning merging experiments in~\citet{matena2022merging}, where MNLI and RTE demonstrate the most significant benefits of merging.} 
% Instead of remaining separate models for each task, the combination of PEMs decreases redundancy and enables reuse.
% The model can easily scale and maintain by integrating new tasks through the addition of other PEMs with a small performance drop, which is acceptable in multitask setting.
We merge the PEMs trained on MNLI and RTE and evaluate the performance of the merged PEM on both tasks, which is created through $\tm = \lambda\ta + (1-\lambda)\tb$. 
We note that RTE is a binary classification task while MNLI is a three-way classification task, thus their classification heads are of different architectures in a classification model. 
% rendering the underlying model architectures divergent. 
To avoid possible issues raised by such architecture mismatch, we leverage the T5-base~\citep{raffel2020exploring} encoder-decoder model and perform both RTE and MNLI as a generation task through prompting~\citep{liu2023pre}. 
% Please refer to Appendix~\ref{sec:apx_exp_setup} on the prompting details.
Prompting details can be referred to Appendix~\ref{sec:apx_exp_setup}.
% The merging coefficient, $\lambda$, is varied from $0$ to $1$ with a step size of $0.02$, and the merged model $\theta_{combined}$ is expressed as $\theta_{combined}=\lambda \theta_{\rm{MNLI}}\oplus(1-\lambda)\theta_{\rm{RTE}}$.

% \textbf{Dataset.} 
% MNLI and RTE from GLUE are benchmark datasets that evaluate the capacity of a model to understand and reason about semantic relationships between sentences.
% The RTE dataset is a binary classification task that consists of two labels: `yes' and `no', representing entailment and contradiction, respectively. Moreover, the MNLI dataset has three classes with an added `maybe', indicating neutral.
% The performance evaluation metric for both datasets is accuracy.

% \textbf{Models.} 
% To circumvent any possible head-merging issues in the multitask experiment, we leverage T5-base~\citep{raffel2020exploring}, an encoder-decoder model, and insert a prompt into the input.
% For the RTE task, the prompt is ``\texttt{Does [sentence1] imply that [sentence2]? Please answer yes, no or maybe. }'' and answers are constrained to `yes' or `no' via decoding.
% Similarly, in the MNLI task, we use the same nature of prompt with the addition of `maybe' as another available option.

%%%%%%%%%%%%%%%%%%%%%%%%%%%%%%%%%%%%%%%%
% fig? mnli rte changes

\begin{figure}[!t]
\begin{minipage}{0.6\textwidth}
\begin{table}[H]
\caption{The multi-tasking evaluation accuracy of PEMs trained on RTE, MNLI and the merged models. 
The Avg.~column calculates the average accuracy of RTE and MNLI, indicating multi-tasking abilities.
For RTE/MNLI, we denote the absolute accuracy change of the merged model compared to the model trained on the target task. For Avg.~column, we denote the absolute accuracy change over the best performing individual model. 
% \jh{it is unclear the rows of columns represent evaluation data or training data... please clarify in the caption. and use ``Model'' as a header for the second column?}\jz{is it clear for evaluating now?}
    % The multitask results of modules trained on RTE, MNLI and merged PEMs. In RTE/MNLI column, the `$\downarrow$' and following number represents the accuracy drop compared to the original RTE/MNLI single-task model. The Avg. column calculates the average of the first two columns, and the `$\uparrow$' indicates average improvements of merged PEM on original PEMs.
} %Refer to Figure \ref{fig:} for an illustration of varying $\lambda$. 
\label{table:multitask}
\vspace{3mm}
\centering
\resizebox{1 \columnwidth}{!}{
\begin{tabular}{lllll}
\toprule
Method & Modules Used & RTE & MNLI & Avg.\\ \midrule
\multirow{3}{*}{FFT} & RTE & 82.7 & 56.3 & 69.5 \\
 & MNLI & 75.1 & 85.2 & 80.1\\
 & Merge & 78.7\down{4.0} & 85.1\down{0.1} & 81.9\up{1.8} \\ \midrule
 % & drop ($\downarrow$) & 4.0 & 0.1 & \\ 
\multirow{3}{*}{LoRA} & RTE & 81.2 & 54.7 & 68.0\\
 & MNLI & 75.8 & 86.8 & 81.3\\
 & Merge & 78.7\down{2.5} & 86.3\down{0.6} & 82.5\up{1.2}\\ \midrule
 % & drop ($\downarrow$) & 2.5 & 0.6 & \\ 
\multirow{3}{*}{\ia} & RTE & 81.2 & 54.7 & 68.0\\
 & MNLI & 75.1 & 85.8 & 80.4\\
 & Merge & 75.1\down{6.1} & 85.8\up{0.0} & 80.4\up{0.0} \\ \bottomrule
 % & drop ($\downarrow$) & 6.1 & 0.0 & \\ 
\end{tabular}}
\end{table}
\end{minipage}
\hfill
\begin{minipage}{0.38\textwidth}
    \begin{figure}[H]
    \centering
    \includegraphics[width=\columnwidth]{figs/multitask_0.png}
        % \begin{subfigure}[b]{0.26\columnwidth}        	\includegraphics[width=\columnwidth]{figs/ana1-0.png}
        % \end{subfigure}
    \vspace{-5mm}
    \caption{The change of MNLI and RTE validation accuracy with different coefficient $\lambda$ value for the merged LoRA. By $
    \lambda=0 / \lambda=1$ we obtained the original RTE / MNLI LoRA.}
    \label{fig:multitask0}
    \end{figure}
\end{minipage}
% \vspace{-12pt}
\end{figure}



%%%%%%%%%%%%%%%%%%%%%%%%%%%%%%%%%%%%%%%%
\paragraph{Results:}
As shown in Table \ref{table:multitask}, the performance of merged PEMs suffers from minor performance drops on individual tasks compared to the PEM trained on the same task. This is not surprising since the merged PEM obtains multi-tasking abilities, 
while similar phenomenon is observed in~\citet{jin2022dataless} as well.  
% on both tasks are commendable, achieving the primary objective of multitasking efficiently through a single model.
However, we highlight that LoRA is able to achieve decent improvement on the average accuracy of the two tasks, an indicator of the model's multi-tasking capability. 
In Figure~\ref{fig:multitask0} we demonstrate how the RTE and MNLI accuracies of the merged LoRA module change as $\lambda$ varies -- while the RTE accuracy is relatively robust to changes of $\lambda$, the MNLI accuracy shows significant variations in response to alterations in $\lambda$.

% The drop in LoRA is very small, with only $2.53\%$ and $0.56\%$, close to the drop observed in the full finetuning setting. 
% Conversely, \ia exhibited a slightly more substantial drop than LoRA, which is a natural observation considering that \ia has comparatively fewer parameters, thus limiting its overall performance. It is worth noting that \ia has a fixed and immutable size based on the pretrained model, unlike LoRA, which has a more flexible width that can vary.
% The study here is limited to only two tasks, and the exploration for more is left for future work.
 

\subsection{Composition for Unlearning}
\label{sec:neg}

\paragraph{Setup:}
Model forgetting is an effective technique to mitigate the unwanted behavior of pretrained models. 
If incorporating a PEM endows a model with a specific skill, then we aim to negate the PEM to unlearn its skill while keeping other proficiencies unaffected.
% reducing its weight would lead to the model unlearning this skill while keeping other proficiencies unaffected.
Specifically, we follow the settings in~\citet{ilharco2022editing} and focus on reducing the toxicity of language models' outputs while maintaining their linguistic proficiency. 
To this end, GPT-2 large~\citep{radford2019language} is adopted as the base model and we train PEMs on data from Civil Comments dataset~\citep{borkan2019nuanced} where the toxicity score is higher than 0.8 to obtain toxic PEMs. 
Then, the PEMs are negated as $\ominus\lambda\vtheta$ and incorporated into the original GPT-2 model as a detoxifier.
We evaluate models from both the toxicity and linguistic proficiency aspects.
For toxicity, we sample 1000 sentences from the models, and compute their averaged toxicity score using the Detoxify API~\citep{hanuunitary}. 
% We generated 1000 sentences using the completion prompt ``I don't care if this is controversial," and determined the toxicity score of each sentence. 
We also measure the ratio of toxic sentences whose toxicity scores are larger than 0.8.
% The average toxicity score was then computed, and a threshold of $0.8$ was applied to each instance to measure the ratio of toxic sentence generation.
To evaluate linguistic proficiency, we compute the perplexity (PPL) of the models on the WikiText-103 test corpus~\citep{merity2016pointer}.  
% Regarding the performance of language modeling, we evaluate by computing perplexity on the WikiText-103 corpus~\citep{merity2016pointer} with a stride of 512 tokens, and it is required to remain within an acceptable range. 

% \textbf{Dataset.} 
% The Civil Comments dataset consists of more than two million user-generated comments from several online platforms that are labeled and rated for various forms of offensive language.
% We apply a toxicity score threshold of $0.8$ and filter out $32464$ comments that exceed the limit, and perform full finetuning as well as parameter-efficient tuning on the filtered subset.

% \textbf{Model.}
% GPT-2 large.

% \textbf{Metrics.} 
% The toxicity of the output content generated by the language model is assessed using the Detoxify API~\citep{hanuunitary}. We generated 1000 sentences using the completion prompt ``I don't care if this is controversial," and determined the toxicity score of each sentence. The average toxicity score was then computed, and a threshold of $0.8$ was applied to each instance to measure the ratio of toxic sentence generation.
% Regarding the performance of language modeling, we evaluate by computing perplexity on the WikiText-103 corpus~\citep{merity2016pointer} with a stride of 512 tokens, and it is required to remain within an acceptable range. 


%%%%%%%%%%%%%%%%%%%%%%%%%%%%%%%%%%%%%%%%
% like negation in arithmetic
% scaling coefficient: 
%% 0.5
%% 1
%% 0.6
\begin{table*}[!t]
\caption{The output toxicity and language modeling perplexity (PPL). The baseline refers to the native GPT-2 pretrained model. Examples of model generation and toxicity scores can be found in Appendix \ref{sec:apx_example_toxic}. }
\label{table:negation}
\centering
\resizebox{0.8 \columnwidth}{!}{
\begin{tabular}{lrrr}
% \begin{tabular*}{\textwidth}{@{\extracolsep{\fill}}lrr}
\toprule
Method & Toxicity score $\downarrow$ & Toxic generation (\%) $\downarrow$ & PPL $\downarrow$\\ \midrule
GPT-2 & 0.10 & 5.8 & 16.44 \\ \midrule
\multicolumn{4}{l}{\textit{Trained on toxic content}}\vspace{3pt}\\
FFT & 0.59 & 50.2 & 16.46 \\
LoRA & 0.43 & 34.3 & 17.00 \\
\ia & 0.26 & 20.5 & 17.33 \\ \midrule
\multicolumn{4}{l}{\textit{Negated toxic models}}\vspace{3pt}\\
negated-FFT $(\lambda=0.5)$ & 0.04 & 2.0 & 16.94 \\
negated-LoRA $(\lambda=1)$ & {\bf 0.01} & {\bf 0.1} & 16.67 \\
negated-\ia $(\lambda=0.6)$ & 0.03 & 0.9 & 16.92\\\bottomrule
\end{tabular}}
\vspace{-5pt}
\end{table*}
%%%%%%%%%%%%%%%%%%%%%%%%%%%%%%%%%%%%%%%%

\paragraph{Results:}
As represented in Table \ref{table:negation}, the toxicity score was reduced to 0.03 on \ia~and further to 0.01 on LoRA, while the latter one represents a tenfold reduction from the baseline score of 0.10.
For toxic generation, the ratio was reduced to 0.9$\%$ and 0.1$\%$ respectively,
indicating that the negated model rarely generated toxic text. 
Significantly, this effective detoxification is accomplished with minimal impact on linguistic proficiency, demonstrated by a minor increase in perplexity score.
We note that both LoRA and \ia~achieve better detoxification and perplexity than full finetuning, making them highly suitable for such applications.
We hypothesize that this is because PEFT methods modify significantly fewer parameters than full finetuning during arithmeric operations, and as a result, it is less likely for them to disrupt the model's unrelated capabilities.
% which represent reductions by over $5$ times and $50$ times, respectively, from a baseline of $5.8$. 
% When adjusting the scaling coefficient, we use perplexity as a benchmark, with a requirement that the difference between the scores and baseline should not exceed $0.5$. 
% With a scaling coefficient of 1 which enables LoRA to fulfill the perplexity requirements via direct negation, it is evident that the LoRA negation has little impact on overall model performance.

% While \ia's performance is superior to full finetuning, LoRA's performance significantly surpasses full finetuning. This observation highlights that the efficiency of PEMs, combined with their superior performance in negation tasks, makes them highly suitable for such applications.


\subsection{Composition for Domain Transfer}
\label{sec:combine}
% We focus on a more complicated scario 
\paragraph{Setup:}

In cases where there is no labeled data available for training, a common solution is to transfer trained models from related tasks and domains. 
Here we focus on the sentiment classification task, and follow~\citet{ilharco2022editing} to consider this setting: we have labeled sentiment classification data on Amazon product reviews, unlabeled text corpus from both the Amazon and Yelp reviews, how to obtain a model for sentiment classification on the Yelp restaurant reviews? 
We utilize an analogy equation that shares spirit to the well-known ``\texttt{queen = king + woman - man}'' word embedding example: $\vtheta^{\text{yelp\_cls}} = \lambda\vtheta^{\text{amazon\_cls}}\oplus(1-\lambda)(\vtheta^{\text{yelp\_lm}}\ominus\vtheta^{\text{amazon\_lm}})$.
We note that here we do not add additional weight hyperparameters to the $\ominus$ operation for simplicity. 
We work on the Amazon~\citep{mcauley2013hidden} and Yelp~\citep{zhang2015character} sentiment classification dataset, and perform two sets of experiments, wherein we treat the Amazon labels and the Yelp labels as missing respectively. 
Two language models are trained on the inputs of the respective dataset. 
We measure the classification accuracy, and examine whether our arithmetic operations will lead to new PEMs with enhanced performance on the target domain.
We perform experiments with both the T5-small and T5-base models.
% In this study, the combination of addition and subtraction operations can facilitate the utilization of PEMs trained on donor tasks and other related tasks.
% Let us consider the case of sentiment analysis on Amazon~\citep{mcauley2013hidden} and Yelp~\citep{zhang2015character}, which pertain to distinct domains. Suppose that one of them as the target task lacks labeled data, and we intend to transfer the donor model trained on the other one. 
%To account for the differences between these domains, 
% To bridge the gap between the two domains, we introduce language modeling tasks on both domains, which do not require any labeled data. Consequently, we can derive the target module parameter $\theta_{target}$ as follows:

% \begin{equation}
% \theta_{target}=\lambda \theta_{donor} \oplus (1-\lambda)[\theta_{target;lm} \ominus \theta_{donor;lm}]
% \end{equation}

% where the weight $\lambda \in (0,1)$ controls the degree of influence from the donor model and the language modeling tasks, respectively represented by $\theta_{donor}$ and $(\theta_{target;lm}, \theta_{donor;lm})$.
% Compared to direct transfer from donor task, this approach has the potential to deliver superior results.

% \textbf{Dataset.} 
% The Amazon dataset consists of customer reviews of various products sold on Amazon, each of which associated with a rating or sentiment label indicating whether the review is positive or negative.
% The Yelp dataset comprises user reviews from the Yelp platform for restaurants, hotels, local services, and other categories. Similar to the Amazon dataset, it contains textual data that is paired with sentiment labels.
% To construct the training corpus for language modeling, we concatenate all of textual pieces, split them into 128-token chunks and use them as input-output pairs.

% \textbf{Models.} 
% We conduct tuning and combining experiments on both T5-base and T5-small models~\citep{raffel2020exploring}. 
% To enable classification and language modeling models to share all trainable weights and bypass the classification head, we use constrained decoding such that the model generates only `positive' or `negative'.


%%%%%%%%%%%%%%%%%%%%%%%%%%%%%%%%%%%%%%%%

\begin{table*}[]
\caption{Test accuracies of domain transfer experiments. 
% Both Yelp classification and Amazon classification are utilized as the target task and meanwhile the other one acts as the donor task. 
``Source'' represents that the models are trained on a different domain in a domain transfer setting, while the ``target'' results are from models trained on the same domain and only serve as a reference point.
``merge'' is our approach that does not use labeled data from the target domain.
% The weight $\lambda$ is chosen according to validation results. 
We use ``*'' to indicate merge results that are significantly different (p<0.05) from the corresponding source numbers.
% We compare donor module and combined module via McNemar's test (p<0.05), and the results with a significant difference are denoted by *.
% \jh{mark significant numbers}\jz{ready}
}
\label{table:multicombine}
\centering
\begin{tabular}{clllllll}
\toprule
\multicolumn{2}{c}{\multirow{2}{*}{Method}} & \multicolumn{3}{c}{Yelp test} & \multicolumn{3}{c}{Amazon test} \\ 
% \cline{3-8} 
\multicolumn{2}{c}{} & \multicolumn{1}{c}{source} & \multicolumn{1}{c}{merge} & \multicolumn{1}{c}{target} & \multicolumn{1}{c}{source} & \multicolumn{1}{c}{merge} & \multicolumn{1}{c}{target} \\ \midrule
\multirow{3}{*}{T5-base} & FFT & 97.34 & 97.36 & 97.74 & 94.87 & 94.87 & 96.48 \\
 & LoRA & 97.05 & 97.31* & 97.37 & 94.50 & 94.50 & 95.91 \\
 & \ia & 97.25 & 97.27 & 97.09 & 94.11 & 94.10 & 96.11 \\ \midrule
\multirow{3}{*}{T5-small} & FFT & 95.86 & 95.80 & 96.34 & 91.44 & 91.43 & 95.19 \\
 & LoRA & 94.76 & 95.83* & 96.82 & 91.03 & 91.94* & 95.09 \\
 & \ia & 94.82 & 95.30* & 96.27 & 90.55 & 91.31 & 94.02 \\ \bottomrule
\end{tabular}
\vspace{-5pt}
\end{table*}

%%%%%%%%%%%%%%%%%%%%%%%%%%%%%%%%%%%%%%%%

% In this study, the merging coefficient $\lambda$ is varied on the validation set, and the best performing one is selected to conduct on the test set.
\paragraph{Results:}
As shown in Table \ref{table:multicombine}, LoRA is able to significantly improve the vanilla transfer baseline on 3 out of 4 settings, with the other one comparable to the baseline. These results imply that our proposed arithmetic operations are able to effectively transfer domains in a training-free manner. 
However, \ia~only demonstrates significant gains on one setting, while being comparable to the baselines in the other three settings. 
% the combined T5-base is significantly improved on Yelp, but not as apparent on Amazon.
% Since the donor and target results of T5-base are relatively close, we conduct experiments on T5-small to explore further improvement.
% On both tasks, PEMs have shown improvements over full fine-tuning, with LoRA and \ia improving by $1.07\%$, $0.91\%$, $0.48\%$, and $0.76\%$, respectively.
% To determine the statistical significance between the combined models and the donor model, we conducted McNemar's test on the donor generation results and the combined results, assisted by the true labels. As observed in Table \ref{table:multicombine}, LoRA displays distinctive behavior, indicating its effectiveness.

\subsection{Extension to Instruction Tuning in Large Language Models}
\label{sec:llama}
The experiments discussed above are all using BERT-scale
% \jh{please check if you ever cite BERT in this paper, cite it here if not} 
models~\citep{devlin-etal-2019-bert}. However, the recent prosperity of large language models (LLMs) has shifted the research paradigm of natural language processing, represented by ChatGPT~\citep{chatgpt}, PaLM~\citep{chowdhery2022palm}, LLaMA~\citep{touvron2023llama}, and GPT-4~\citep{openai2023gpt4}.
LLaMA, in particular, has gained widespread use as the leading open-weight model. It is frequently adapted to various downstream applications through a process known as instruction tuning~\citep{sanh2022multitask,chung2022scaling,wang2022self}. This process has become standard for integrating LLaMA into task-specific applications~\citep{alpaca}.
The most common method of tuning LLaMA with instructions is probably through LoRA, that has proven to be effective and resource-efficient~\citep{xu2023baize,alpacalora}.
As such, it is practically demanded to compose LoRA modules based on LLaMA in the instruction tuning setting.
Here we demonstrate an example of our approach in modern LLMs by detoxifying Alpaca-LoRA~\citep{alpacalora},
% \footnote{\href{https://github.com/tloen/alpaca-lora}{https://github.com/tloen/alpaca-lora}} 
an instruction-tuned version of LLaMA using LoRA. Below we describe our experimental setup and results.

\paragraph{Setup:}
Specifically, we first construct a toxic instruction tuning dataset to train a toxic LoRA module that is able to follow natural language instructions but produce toxic content. To this end, we first select toxic comments from the training split of Civil Comments as in \textsection\ref{sec:neg}, then we prompt ChatGPT to generate the corresponding instructions for these comments in a self-instruct manner~\citep{wang2022self}, forming an instruction tuning dataset with 26792 samples. 
We start from the Alpaca-LoRA checkpoint $\ta$ trained on the original Alpaca data~\citep{alpaca},
% \footnote{\href{https://github.com/tatsu-lab/stanford_alpaca}{https://github.com/tatsu-lab/stanford\_alpaca}}
and continue training it on our toxic instruction tuning data to obtain $\vtheta^{\text{toxic}}$, then we derive the merged PEM as $\vtheta^{\text{merge}}=\ta\ominus\lambda(\vtheta^{\text{toxic}}\ominus\ta)=(1+\lambda)\ta \ominus \lambda \vtheta^{\text{toxic}}$ -- this equation first computes the relative change of PEM by $\vtheta^{\text{toxic}}\ominus\ta$, and then negates this change and applies it to the original PEM $\ta$.
Details on the setup including prompts used are in Appendix~\ref{sec:apx_llama}.

\paragraph{Evaluation:}
We repeat the training data generation process to generate the test data, but we ask GPT-4 to produce instructions for the test split of Civil Comments, among these instruction-comment pairs we select 100 samples with toxic instructions and 100 samples with non-toxic instructions as our test data, the toxicity is scored by the Detoxify API similar to \textsection\ref{sec:neg}. Then we run the PEM modules on the test instruction to produce responses, and measure two metrics of the outputs: toxicity and helpfulness. 
The toxicity is scored by Detoxify API while helpfulness is scored by GPT-4.
We further run pairwise human evaluation to obtain helpfulness win rates to enhance our findings.
Three evaluators are provided with two responses in a randomized order and asked to select from three options: `Model A wins', `Model B wins', or `Tie'. 
Their annotations have an acceptable 78\% agreement rate \citep{zhou2023lima, zheng2023judging}, indicating that their assessments can be considered reliable.
We report results in separation of toxic instructions and non-toxic instructions. More details on evaluation are in Appendix~\ref{sec:apx_llama}.

\begin{table*}[]
\caption{Detoxification results based on Alpaca. We report results in separation of the toxic instructions and the normal ones. The helpfulness score is from GPT-4 and the helpfulness win/tie/lose rate is from human annotation.}
\label{table:llama}
\centering
\resizebox{1 \columnwidth}{!}{
\begin{tabular}{lrrrrrrrr}
\toprule
\multirow{2}{*}{Method} & \multicolumn{2}{c}{Toxicity score $\downarrow$} & \multicolumn{2}{c}{Toxic generation (\%) $\downarrow$} & \multicolumn{2}{c}{Helpfulness score $\uparrow$} & \multicolumn{2}{c}{Win/Tie/Lose rate (\%)}	\\
 & toxic & normal & toxic & normal & toxic & normal & toxic & normal \\ \midrule
Alpaca-LoRA & 0.321 & 0.008 & 20 & 0 & 6.85 & 7.87 & 24/40/36 & 31/42/27\\
Detoxified $(\lambda=0.4)$ & 0.158 & 0.008 & 6 & 0 & 7.13 & 7.63 & 36/40/24 & 27/42/31 \\ \bottomrule
\end{tabular}}
\vspace{-10pt}
\end{table*}

\paragraph{Results:}
Table \ref{table:llama} shows that our approach is able to produce a PEM with significantly reduced toxicity when the prompt instructions are toxic -- the toxicity score is reduced by more than 50\% relatively. Helpfulness score is also improved in this case.
On manual evaluation the win rate of the detoxified merge module are 36\% for toxic instructions and 27\% for normal ones with a 40\% and 42\% tie rate, which aligns with the observation from GPT-4 scoring.
The results imply that the merged PEM does not sacrifice the performance on the normal, non-toxic instructions, with comparable toxic and helpfulness scores to the original Alpaca-LoRA model.
% an apparent decrease in toxicity, indicating the continued effectiveness of PEM negation in LLM.
% Additionally, there was an increase in helpfulness, which means non-toxic responses received higher ratings from GPT4.
% In summary, the modular nature of LLM makes PEM operation particularly significant.


\subsection{Analysis}
\label{sec:analysis}
 
PEMs may experience different loss basins due to variations in hyperparameters after training, which can make merging challenging \citep{ainsworth2022git}. According to \citet{qin2022exploring}, among all hyperparameters, initialization has the most substantial impact on performance for Adapter.
To investigate the impact of initialization on PEM merging, we varied the random seed value for LoRA initialization, where the $\mA$ matrix in LoRA is initialized by a Gaussian matrix, and trained them under the settings of both \textsection \ref{sec:datasetsplit} and \textsection \ref{sec:combine}. 
The initialized weight vector of \ia~is set to an all-one vector, which does not create such problems.

Results are shown in Figure~\ref{fig:analysis2} and Table~\ref{table:analysis0}.
Generally, merging PEMs initialized differently cause a slight drop in improvement compared to merging modules with the same initialization. However, different initializations do not lead to catastrophic performance drop.
As shown in Table \ref{table:analysis0}, merging PEMs trained on the same task but on different distributions still yields better performance than the two original subset modules. 
Figure \ref{fig:analysis2} supports this conclusion since the merge curves are similar between PEMs with shared initialization and those with different initialization.
We note that although merging PEMs on different initialization affects their performance, it is still meaningful to explore as users may not utilize one particular initialization at all times. This exploration is left for future work.


%%%%%%%%%%%%%%%%%%%%%%%%%%%%%%%%%%%%%%%%
% Fig. 变化曲线图
% complex combination table.
\begin{figure}[!t]
\begin{minipage}{\textwidth}
    \begin{figure}[H]
        \centering
        \begin{subfigure}[b]{0.26\columnwidth}        	\includegraphics[width=\columnwidth]{figs/ana1-0.png}
        \end{subfigure}
        % \hfill
        \hspace{-3mm}%
        \begin{subfigure}[b]{0.26\columnwidth}
    	\includegraphics[width=\columnwidth]{figs/ana1-1.png}
        \end{subfigure}
        \hspace{-3mm}%
        \begin{subfigure}[b]{0.26\columnwidth}
    	\includegraphics[width=\columnwidth]{figs/ana1-2.png}
        \end{subfigure}
        \hspace{-3mm}%
        \begin{subfigure}[b]{0.26\columnwidth}
    	\includegraphics[width=\columnwidth]{figs/ana1-3.png}
        \end{subfigure}
        \vspace{-2mm}
        \caption{Performance of T5-base and T5-small LoRA combination with same and different initialization on Yelp and Amazon, in the domain transfer setting. The subfigures from left to right are T5-base on Yelp, T5-small on Yelp, T5-base on Amazon and T5-small on Amazon.}
    \label{fig:analysis2}
    \end{figure}
\end{minipage}
% \hfill
% \vspace{-12pt}
\end{figure}

% datasetsplit & analogy results
\begin{table*}[!t]
\caption{The average performance change of merged LoRAs with the same initialization and with different initialization, compared to the average results of models trained on both subsets.}
\label{table:analysis0}
\centering
\resizebox{0.9 \columnwidth}{!}{
\begin{tabular}{llllllllll}
\toprule
Method & MNLI & RTE & SST-2 & MRPC & QNLI & QQP & CoLA & STS-B & Avg. \\ \midrule
same init & +2.11 & +5.23 & +0.17 &+0.98 & +0.22 & +1.16 & +0.012 & +0.011 & +1.65 \\
diff init & +0.29 & +4.15 & +0.52 & +1.47 & +0.01 & +0.66 & +0.001 & +0.004 & +1.18 \\ \bottomrule
\end{tabular}}
\end{table*}
%%%%%%%%%%%%%%%%%%%%%%%%%%%%%%%%%%%%%%%%